%=============================================================================
% SECTION 8C: CONVENTION-LOCKED α FROM THE MICROSECTOR
% Part V: Parameter-Free Predictions (continued)
% The microsector derivation with forced fork and no hidden knobs
%=============================================================================

\section{Convention-Locked \texorpdfstring{$\alpha$}{α} from the Microsector}\label{sec:alpha-locked}

The preceding sections derived $\alpha$-relations from gauge emergence and electroweak physics. This section presents the \textit{microsector completion}: a derivation of $\alpha^{-1} = 137.036$ from the internal geometry~\cite{AlcockAbInitioAlpha2026}, with all conventions locked and no hidden tuning parameters. The result matches experiment at sub-ppm precision.

%-----------------------------------------------------------------------------
\subsection{Design Constraint: No Hidden Tuning Parameters}\label{sec:no-knobs-policy}
%-----------------------------------------------------------------------------

We impose a \emph{no-knobs policy}: once the microsector geometry, bundle data, and truncation level are fixed, the predicted $\alpha$ must be stable without invoking subleading heat-kernel terms as ppm-level tuners. Concretely, we choose a cutoff rule that prevents $a_6, a_8, \ldots$ from acting as free correction dials (Sec.~\ref{sec:alpha-regularization}).

\paragraph{Motivation.}
Any theory that ``predicts'' a fundamental constant but allows ppm-level adjustments via regulator moments or trace normalizations is not truly predictive---it has hidden knobs. The microsector completion must lock all such freedoms.

%-----------------------------------------------------------------------------
\subsection{Operator Choice (Locked)}\label{sec:alpha-operator-choice}
%-----------------------------------------------------------------------------

On the internal microsector $X = \mathbb{C}P^2 \times S^3$, we take a Laplace-type operator given by the \emph{connection Laplacian}:
\begin{equation}\label{eq:connection-laplacian}
P = -g^{ij}\nabla_i\nabla_j,
\end{equation}
acting on the internal bundle that carries the emergent gauge degrees of freedom.

\paragraph{Bundle structure.}
The U(1) factor is implemented via twisting by a line bundle over $\mathbb{C}P^2$ with curvature proportional to the K\"ahler form $\omega$, taken trivial over $S^3$. This choice is minimal and convention-stable: the K\"ahler form is parallel ($\nabla\omega = 0$), so derivative terms in higher Seeley--DeWitt coefficients vanish automatically.

\paragraph{Why this is locked.}
The gauge-kinetic extraction from $a_4$ is unambiguous with this operator choice. Alternative operators would introduce additional terms proportional to curvature derivatives, creating ppm-level ambiguities. The connection Laplacian with parallel curvature eliminates this freedom.

%-----------------------------------------------------------------------------
\subsection{Regularization/Truncation Rule (Locked)}\label{sec:alpha-regularization}
%-----------------------------------------------------------------------------

We define the spectral action with a \emph{plateau cutoff} function $f$:
\begin{equation}\label{eq:spectral-action-plateau}
S = \mathrm{Tr}\, f(P/\Lambda^2),
\end{equation}
where $f$ is constant in a neighborhood of the origin.

\paragraph{The plateau condition.}
Equivalently, $f^{(n)}(0) = 0$ for all $n \geq 1$, so all negative moments vanish:
\begin{equation}
f_{-2} = f_{-4} = \cdots = 0.
\end{equation}

\paragraph{Why this is locked.}
This eliminates the possibility of using $a_6$ (or higher) contributions as hidden ppm-level tuning knobs. With generic smooth cutoffs (e.g., Gaussian), the $a_6$ contribution would be $\sim 2\%$---far too large and requiring fine-tuned cancellation. The plateau cutoff is the unique choice that:
\begin{enumerate}
\item Preserves the leading $a_4$ gauge kinetic term
\item Eliminates subleading heat-kernel contributions
\item Requires no moment-tuning
\end{enumerate}

%-----------------------------------------------------------------------------
\subsection{Finite-$k$ Truncation and the $(k+3)/(k+4)$ Factor (Locked)}\label{sec:alpha-finite-k}
%-----------------------------------------------------------------------------

We implement a finite-$k$ truncation via Toeplitz quantization at level $m = k + 3$ on $\mathbb{C}P^1$, where:
\begin{equation}\label{eq:d-dimension}
d = \dim H^0(\mathbb{C}P^1, \mathcal{O}(m)) = m + 1 = k + 4.
\end{equation}

\paragraph{Origin of the +3 shift.}
The shift $m = k + 3$ arises from the Spin$^c$ structure on $\mathbb{C}P^2$:
\begin{equation}
K_{\mathbb{C}P^2} = \mathcal{O}(-3) \quad \Rightarrow \quad L_{\det} = K^{-1} = \mathcal{O}(3).
\end{equation}
When restricting to $\mathbb{C}P^1 \subset \mathbb{C}P^2$, the line bundle $\mathcal{O}(k) \otimes L_{\det}$ becomes $\mathcal{O}(k+3)$, giving sections of dimension $k + 4$.

\paragraph{The spectral cutoff.}
The determinant-channel removal at finite $d$ fixes the spectral cutoff as:
\begin{equation}\label{eq:lambda-cutoff}
\Lambda^3 = k \cdot \frac{d-1}{d} = k \cdot \frac{k+3}{k+4}.
\end{equation}

This is the unique finite-size factor permitted by the truncation rule; it is \emph{not} inserted to improve agreement.

%-----------------------------------------------------------------------------
\subsection{The Forced Microsector Fork}\label{sec:alpha-forced-fork}
%-----------------------------------------------------------------------------

At this point there is a \emph{forced binary fork}, determined solely by what finite Hilbert space carries the microsector trace.

\subsubsection{Branch A: Regular-Module Microsector (Survives)}

Take the finite Hilbert space to be the algebra itself:
\begin{equation}\label{eq:HF-regular}
\mathcal{H}_F := A = M_d(\mathbb{C}),
\end{equation}
with Hilbert--Schmidt inner product $\langle X, Y \rangle = \mathrm{Tr}(X^\dagger Y)$, and gauge action by inner derivations:
\begin{equation}
\mathrm{ad}_a(X) = [a, X].
\end{equation}

\paragraph{Trace normalization.}
The UV-normalized trace is naturally the democratic normalization per matrix degree of freedom:
\begin{equation}\label{eq:tr-dem}
\mathrm{tr}_{\rm dem}(\cdot) := \frac{1}{d^2}\,\mathrm{Tr}_{\mathcal{H}_F}(\cdot).
\end{equation}

\paragraph{Conversion to physics normalization.}
When reporting the final gauge kinetic term in canonical generator normalization on $\mathfrak{su}(d)$:
\begin{equation}
\mathrm{tr}_{\mathfrak{su}}(\cdot) = \frac{1}{d^2 - 1}\,\mathrm{Tr}_{\mathfrak{su}}(\cdot),
\end{equation}
the conversion factor is \emph{forced}:
\begin{equation}\label{eq:boost-factor}
\boxed{\varepsilon_{\rm adj}^{(A)} = \frac{d^2}{d^2 - 1}}
\end{equation}

For $k = 60$, $d = 64$:
\begin{equation}
\varepsilon_{\rm adj}^{(A)} = \frac{4096}{4095} = 1.000244\ldots
\end{equation}

\subsubsection{Branch B: Fermion-Representation Microsector (Falsified)}

If instead the kinetic term trace is taken over a $d$-dimensional fermion representation space $\mathcal{H}_F \cong \mathbb{C}^d$ (as in conventional matter spectral triples), unimodularity literally removes the identity generator channel, yielding the drop factor:
\begin{equation}\label{eq:drop-factor}
\varepsilon_{\rm adj}^{(B)} = \frac{d^2 - 1}{d^2} = \frac{4095}{4096} = 0.999756\ldots
\end{equation}

%-----------------------------------------------------------------------------
\subsection{Decision Rule and Lock}\label{sec:alpha-decision-rule}
%-----------------------------------------------------------------------------

Holding \emph{all other ingredients fixed} (geometry, $g_F$, hypercharge trace, and the finite-$k$ rule $\Lambda^3 = k(k+3)/(k+4)$), we compute $\alpha^{-1}$ under both microsector trace choices.

\paragraph{Numerical results.}

\begin{table}[h]
\centering
\caption{Microsector fork: numerical comparison at $k = 60$.}\label{tab:microsector-fork}
\renewcommand{\arraystretch}{1.3}
\begin{tabular}{lccc}
\toprule
Branch & Factor & $\alpha^{-1}$ & Residual (ppm) \\
\midrule
\textbf{A (regular-module)} & $\displaystyle\frac{4096}{4095}$ & $137.03599985$ & $-0.006$ \\[8pt]
B (fermion-rep) & $\displaystyle\frac{4095}{4096}$ & $137.03014445$ & $+42.7$ \\
\midrule
Experimental & --- & $137.035999084$ & --- \\
\bottomrule
\end{tabular}
\end{table}

\paragraph{Branch A: matches.}
The regular-module microsector matches $\alpha^{-1}$ at sub-ppm level \emph{without} invoking higher heat-kernel terms (consistent with the plateau cutoff).

\paragraph{Branch B: cannot be rescued.}
The fermion-rep microsector misses by $\sim 43$ ppm. This deficit \textbf{cannot} be repaired by:
\begin{itemize}
\item Adjusting the U(1)/non-Abelian mixing weights ($w$): would require $\Delta w/w = -200\%$
\item Adjusting $g_F$: would require $\Delta g_F/g_F = +200\%$
\item Using $a_6$ correction: would require tuning cutoff moments to $f_{-2}/f_0 \sim 10^{-3}$, violating the no-knobs policy
\end{itemize}
\paragraph{The lock.}\mbox{}

\begin{tcolorbox}[colback=green!5!white, colframe=green!50!black, title=Microsector Lock]
Under the no-knobs policy, we \textbf{adopt} the regular-module microsector completion (Branch A) and treat Branch B as \textbf{falsified}.

\textbf{Committed microsector:}
\begin{itemize}
\item Hilbert space: $\mathcal{H}_F = A = M_d(\mathbb{C})$ (regular module)
\item Dimension: $\dim(\mathcal{H}_F) = d^2 = 4096$
\item Gauge action: inner derivations $\mathrm{ad}_a(X) = [a, X]$
\item UV trace: $\mathrm{tr}_{\rm dem} = (1/d^2)\,\mathrm{Tr}$
\item Factor: BOOST $= d^2/(d^2 - 1) = 4096/4095$
\end{itemize}
\end{tcolorbox}

%-----------------------------------------------------------------------------
\subsection{The Complete Derivation Chain}\label{sec:alpha-chain}
%-----------------------------------------------------------------------------

The $\alpha$ derivation is now fully locked:

\begin{table}[h]
\centering
\caption{Complete derivation chain for $\alpha^{-1}$.}\label{tab:alpha-chain}
\renewcommand{\arraystretch}{1.2}
\resizebox{\columnwidth}{!}{%
\begin{tabular}{llll}
\toprule
Component & Value & Source & Status \\
\midrule
$K_{\mathbb{C}P^2} = \mathcal{O}(-3)$ & $-3$ & Algebraic geometry theorem & Rigorous \\
$L_{\det} = K^{-1}$ & $\mathcal{O}(3)$ & Spin$^c$ structure & Rigorous \\
$d = k + 4$ & $64$ & $\dim H^0(\mathcal{O}(k+3))$ & Rigorous \\
$(d-1)/d$ & $63/64$ & Traceless projection & Derived \\
$N_{\rm species}$ & $7$ & SM SU(2) components & SM content \\
$\mathrm{Tr}(Y^2)$ & $10$ & SM hypercharges & SM content \\
$g_F$ & $8$ & Spectral triple ($J \times \gamma \times \mathbb{C}$) & Derived \\
$w = N_{\rm species}/(g_F \cdot \mathrm{Tr}(Y^2))$ & $7/80$ & Hypercharge weighting & Derived \\
$\varepsilon_{\rm adj}$ & $4096/4095$ & Regular-module trace conversion & Forced \\
\midrule
$\alpha^{-1}$ & $137.03599985$ & All above combined & $< 0.01$ ppm \\
\bottomrule
\end{tabular}%
}
\end{table}

\paragraph{Closure of $k_{\max} = 60$.}
The baseline normalization $\Lambda^3 = 885.9375$ (from $k = 60$, $a = 9$, $n = 5$, $N = 3$) sets the overall scale.  Within the finite-symmetry closure framework adopted in this section, the value $k_{\max}=60$ follows from the following auxiliary structural postulates:
\begin{enumerate}
\item The microsector channel symmetry $G$ acts faithfully on a real three-dimensional generation space.
\item $G$ is orientation-preserving and simple (no hidden normal subgroup).
\item The channel algebra furnishes exactly five conjugacy classes, matching the five chiral multiplet types in one SM generation.
\item Choose the \emph{minimal} such group.
\end{enumerate}
Under these auxiliary postulates, the unique solution is the icosahedral rotation group $G \cong A_5$, hence $k_{\max} = |A_5| = 60$.  This is a \emph{conditional closure theorem inside the finite-symmetry framework}.  It should not be read as a derivation from the core DFD field equation alone.  Its value is that it removes arbitrary integer freedom once the stated structural postulates are adopted.  The independent Bridge Lemma (Appendix~\ref{app:bridge}), lattice Monte Carlo selection (Appendix~\ref{app:kmax-discovery}), and minimal-padding argument then function as nontrivial consistency checks rather than as hidden tuners.  Once $k_{\max}$ is fixed, only discrete choices remain.

\paragraph{Unconditional content.}
What does \emph{not} depend on the auxiliary postulates is the following: once any integer $k_{\max}$ is fixed, the entire microsector output ($\alpha^{-1}$, fermion masses, CKM structure, neutrino spectrum) follows with zero continuous free parameters.  The structural postulates above select $k_{\max}=60$ from the integers; the theory's numerical output is then falsifiable against ${>}30$ independent measurements.

%-----------------------------------------------------------------------------
\subsection{Sharp Falsifier}\label{sec:alpha-falsifier}
%-----------------------------------------------------------------------------

The microsector choice $\mathcal{H}_F = A$ is a \textbf{testable ontological claim}:

\begin{quote}
\emph{``The finite Hilbert space of the DFD Toeplitz microsector is the algebra itself ($M_d(\mathbb{C})$), not a fermion representation space ($\mathbb{C}^d$).''}
\end{quote}

\paragraph{If future work derives $\mathcal{H}_F = \mathbb{C}^d$ from first principles:}
\begin{itemize}
\item DFD fails by 43 ppm
\item Cannot be rescued without fine-tuning
\item Theory requires fundamental revision
\end{itemize}

\paragraph{If future work derives $\mathcal{H}_F = A$ from first principles:}
\begin{itemize}
\item DFD is confirmed
\item BOOST factor is forced, not fitted
\item The $\alpha$ match is genuine
\end{itemize}

%-----------------------------------------------------------------------------
\subsection{The Closed-Form Result}\label{sec:alpha-one-liner}
%-----------------------------------------------------------------------------

Collecting all locked ingredients, the fine-structure constant is given
by a single equation with no continuous free parameters:
\begin{equation}\label{eq:alpha-one-liner}
\boxed{\;
\begin{aligned}
\alpha^{-1} &= \frac{\pi^{3/2}}{24}\,\mathrm{Tr}(Y^2)\,k_{\max}\,
\frac{k_{\max}\!+\!3}{k_{\max}\!+\!4}\\[4pt]
&\quad\times\!\left[1 + \frac{7}{80\!\cdot\!4095}\right]
= 137.036
\end{aligned}
\;}
\end{equation}
where:
\begin{itemize}[nosep]
\item $\pi^{3/2}/24$: geometry factor from the $a_4$ Seeley--DeWitt
  coefficient on $\mathbb{C}P^2 \times S^3$
\item $\mathrm{Tr}(Y^2) = 10$: Standard Model hypercharge trace
  (3 generations of $Q_L, u_R, d_R, L_L, e_R$)
\item $k_{\max} = 60$: topological cutoff from the Bridge Lemma
  (Spin$^c$ index on $\mathbb{C}P^2$, $= |A_5|$)
\item $(k_{\max}+3)/(k_{\max}+4) = 63/64$: Toeplitz truncation
  from the Spin$^c$ determinant line $L_{\det} = \mathcal{O}(3)$
\item $[1 + 7/(80 \times 4095)]$: regular-module microsector
  correction ($4095 = 64^2 - 1 = d^2 - 1$)
\end{itemize}
The exact numerical evaluation via the full Chern--Simons weight sum
gives $\alpha^{-1} = 137.03599985$ (residual $-0.006$~ppm vs.\
experiment).

%-----------------------------------------------------------------------------
\subsection{Summary}\label{sec:alpha-locked-summary}
%-----------------------------------------------------------------------------

\begin{tcolorbox}[colback=blue!5!white, colframe=blue!50!black, title=Summary: Convention-Locked $\alpha$]
\textbf{Result:}
\begin{equation}
\alpha^{-1} = 137.03599985 \quad \text{(residual: } {-}0.006\text{ ppm)}
\end{equation}

\textbf{Locked conventions:}
\begin{itemize}
\item Operator: connection Laplacian with parallel curvature
\item Regulator: plateau cutoff ($f_{-2} = f_{-4} = \cdots = 0$)
\item Finite-$k$: Toeplitz truncation with $d = k + 4 = 64$
\item Microsector: regular-module ($\mathcal{H}_F = M_d(\mathbb{C})$)
\item Trace: democratic UV $\to$ per-generator physics (BOOST forced)
\end{itemize}

\textbf{The fermion-rep microsector is falsified:}
\begin{itemize}
\item 43 ppm deficit cannot be filled
\item All salvage paths blocked (w, $g_F$, $a_6$)
\item Under no-knobs policy, only Branch A survives
\end{itemize}

\textbf{Falsification criterion:} If $\mathcal{H}_F = \mathbb{C}^d$ is derived from microsector first principles, DFD's $\alpha$ prediction fails.
\end{tcolorbox}
